\section{The one-chart calculation}
\label{app:one-chart}

This appendix expands the proof of Theorem~\ref{thm:tailwise-derived} at the
three places where it is compressed: the affine verification of the derived
equivalence, the \(2\)-commutative obstruction-theory square, and the
\(\mu_k\)-equivariance and rigidification step.  The hypotheses, the statement,
and the conclusions of that theorem are unchanged; what follows adds no new
assumption and asserts nothing stronger.

One convention must be fixed before any comparison.  Woodward works with
classical stacks carrying perfect obstruction theories in the sense of
Behrend--Fantechi, over Artin stacks of prestable domains, and with Vistoli's
bivariant Chow groups for representable morphisms; his index complex is
\((Rp_*e^*T(X/G))^\vee\), stated absolutely in
\cite[Definition~7.13(a)]{WoodwardQKIII} and relative over the twisted domain
stack in \cite[Section~8.3]{WoodwardQKIII} and
\cite[Example~6.6(c)]{WoodwardQKII}.  Derived stacks play no role there.  The
derived fixed-section stacks used in Theorem~\ref{thm:tailwise-derived} are an
auxiliary layer of the present manuscript.  Their only purpose is to make the
degree-shift comparison functorial in nilpotent and simplicial thickenings of
the test base, which is what upgrades an isomorphism of perfect complexes to an
isomorphism of obstruction theories \emph{as morphisms to a cotangent complex}.
Every comparison below is returned to Woodward's setting by
Proposition~\ref{prop:app-square}: the classical truncation of the derived
obstruction theory is his relative perfect obstruction theory over the fixed
moduli of marked domains, so the identification is an identification of the
relative virtual classes he uses.

The imported statements are: the clutching description of the rotation-fixed
gauged maps \cite[Definition~9.7, Lemmas~9.8 and~9.9,
Corollary~9.10]{WoodwardQKIII}; the obstruction theory
\cite[Definition~7.13]{WoodwardQKIII} in its relative form
\cite[Section~8.3]{WoodwardQKIII}, \cite[Example~6.6(c)]{WoodwardQKII}; the
cutting axiom \cite[Proposition~7.14(b)]{WoodwardQKIII} together with the fibre
product over the diagonal on which its proof rests
\cite[Proposition~5.21]{WoodwardQKII}; and the rigidified inertia conventions
\cite[Section~4.3, Proposition~4.3(f),(g)]{WoodwardQKII}.  Two standard
comparisons are used and not reproved: the cotangent complex of a derived
mapping stack is the pushforward of the pulled-back cotangent complex of the
target, and a classical obstruction theory built from the deformation theory of
a universal object is the classical truncation of the derived one.
Remark~\ref{rem:app-imports} says exactly where each enters.

\subsection{Graded setup and the derived fixed-section stack}
\label{app:setup}

Throughout, \(W\) is the smooth projective \(\Gm\)-variety of
Definition~\ref{def:gauged-admissible} and we work over \(\C\).  Fix the
generic semistable locus \(W^{\mathrm{gen}}\subset W\) and the two fixed
components \(F_-,F_+\) prescribed by the rotation-fixed graph type, as in the
proof of Theorem~\ref{thm:tailwise-derived}.  The gauge group is \(\Gm\); the
domain rotation group is written \(\C^\times_{\mathrm{rot}}\).

The domain of the principal component is the parametrized line \(P\) with its
two marked points \(0,\infty\) and the two affine charts
\[
 P^0=\operatorname{Spec}\C[z],\qquad
 P^\infty=\operatorname{Spec}\C[z'],\qquad z'=z^{-1},
\]
glued along the open orbit \(\Gm\subset P\).  Each chart carries the rotation
action scaling its own coordinate.  We record an exponent \(a_\pm\in\mathbf Z\)
in each chart's own coordinate, so that both extension conditions below are
conditions at the vanishing locus of the local coordinate; this matches the
convention \(W_F^a=\{x\in W^{\mathrm{gen}}:\lim_{t\to0}t^ax\in F\}\) used in
the proof of Theorem~\ref{thm:tailwise-derived}.

\begin{definition}[Derived fixed-section stack]
\label{def:app-fixed-section}
Let \(R\) be a simplicial commutative \(\C\)-algebra.  An \(R\)-point of
\(\mathfrak S_{a_-,a_+}\) is a rotation-equivariant map
\(u\colon P_R\to[W/\Gm]\) whose rotation-equivariant structure on the
underlying \(\Gm\)-bundle is the clutching datum of
\cite[Definition~9.7]{WoodwardQKIII} for the pair of cocharacters
\(t\mapsto t^{a_\pm}\), together with the prescribed stabilizer lift and graph
component.  Morphisms of \(R\)-points are isomorphisms of these data.  This is
a functor on simplicial commutative \(\C\)-algebras, defined degreewise in the
simplicial direction; characteristic zero is used throughout, both for
simplicial commutative algebras to model connective commutative differential
graded algebras and for exactness of invariants under the reductive groups that
appear.
\end{definition}

By Sumihiro's theorem \(W\) is covered by \(\Gm\)-invariant affine opens.  Fix
one, \(\operatorname{Spec}A\), and write its coordinate ring as a
\(\Gm\)-representation,
\begin{equation}\label{eq:app-grading}
 A=\bigoplus_{w\in\mathbf Z}A_w,
 \qquad
 f\in A_w\iff f(t\cdot x)=t^wf(x).
\end{equation}
Over a test algebra \(R\), the trivialized chart at \(0\) has coordinate ring
\(R[z]\), graded with \(z\) in degree \(1\) and \(R\) in degree \(0\); the open
orbit has \(R[z^{\pm1}]\) with the same grading.

\begin{remark}[What the exponent does and does not change]
\label{rem:app-exponent}
The exponent \(a\) enters the chart only through the twist of the
rotation-equivariant structure on the \(\Gm\)-bundle.  It does not change the
domain, the target, or the trivializations.  This is the reason the
calculations below separate cleanly into a part that depends on \(a\) only
through its sign and a part---the universal section itself---that depends on
\(a\) exactly.
\end{remark}

\subsection{Evaluation at one and the attractor algebra}
\label{app:attractor}

\begin{lemma}[Unique extension on one chart]
\label{lem:app-graded-extension}
Let \(R\) be a simplicial commutative \(\C\)-algebra and \(a\in\mathbf Z\).
\begin{enumerate}[label=\textup{(\alph*)}]
\item Ring maps \(\psi\colon A\to R[z^{\pm1}]\) sending \(A_w\) into the
\(z\)-degree \(aw\) part --- the equivariance condition for a rotation-fixed
section with clutching exponent \(a\), for the gradings
\eqref{eq:app-grading} --- correspond bijectively, degreewise in the
simplicial direction, to ring maps \(x\colon A\to R\), by
\begin{equation}\label{eq:app-graded-map}
 \psi(f)=x(f)\,z^{aw}\qquad(f\in A_w).
\end{equation}
The correspondence is evaluation at \(z=1\).
\item Such a \(\psi\) takes values in \(R[z]\) if and only if \(x\) annihilates
every \(A_w\) with \(aw<0\), that is, if and only if \(x\) factors through the
\emph{attractor algebra}
\begin{equation}\label{eq:app-attractor-ideal}
 A^a:=A/I_a,\qquad I_a:=\bigl(A_w:aw<0\bigr).
\end{equation}
\end{enumerate}
\end{lemma}

\begin{proof}
The grading forces \(\psi(A_w)\subset(R[z^{\pm1}])_{aw}=R\,z^{aw}\), so
\(\psi\) is determined by the \(R\)-linear maps \(x_w\colon A_w\to R\) defined
by \(\psi(f)=x_w(f)z^{aw}\).  Since \(z^{aw}z^{aw'}=z^{a(w+w')}\), the
collection \((x_w)\) is multiplicative if and only if the single map
\(x=\bigoplus_wx_w\colon A\to R\) is a ring map; unitality corresponds to
\(x(1)=1\).  This is (a), and both directions are levelwise constructions in
the simplicial variable, so they hold degreewise.  For (b), \(R[z]\) is the
part of \(R[z^{\pm1}]\) of nonnegative degree, so \(\psi(A_w)\subset R[z]\)
fails exactly when \(aw<0\) and \(x_w\neq0\).  The elements with \(aw<0\)
generate \(I_a\), and a ring map killing a generating set of an ideal kills the
ideal.
\end{proof}

Nothing in Lemma~\ref{lem:app-graded-extension} required \(A\) to be resolved,
and the graded pieces were used only to split a map of sets.  What makes the
levelwise functor \(R\mapsto\operatorname{Hom}(A,R)\) the derived mapping
space is smoothness of the \emph{source} algebra \(A\), which holds because
\(W\) is smooth; the target \(R[z]\) plays no role in that point.  The next two statements record the
consequences that the proof of Theorem~\ref{thm:tailwise-derived} uses.

\begin{lemma}[Sign dependence]
\label{lem:app-sign}
For \(a\neq0\) the ideal \(I_a\) of \eqref{eq:app-attractor-ideal} depends only
on the sign of \(a\).  Consequently, for two exponents \(a,a'\) of the same
sign, the attractor algebras \(A^a=A^{a'}\), the induced strata, and the
restriction maps to the open orbit coincide; only the universal section
\(z\mapsto z^ax\) of \eqref{eq:app-graded-map} differs.  For \(a=0\) one has
\(I_0=0\) and \(A^0=A\).
\end{lemma}

\begin{proof}
The condition \(aw<0\) is the condition \(\operatorname{sign}(a)w<0\), which
depends on \(a\) only through its sign; for \(a=0\) it is empty.  The strata
and the restriction maps are defined from \(A^a\) and the inclusion
\(\operatorname{Spec}A^a\subset\operatorname{Spec}A\), so they inherit this.
The section \eqref{eq:app-graded-map} is \(z\mapsto z^ax\), which changes with
\(a\).
\end{proof}

The value \(a=0\) is the degenerate case in which the extension condition is
vacuous, the limit point need not be fixed, and no fixed component \(F\) is
selected.  This is why exponent zero occurs in
Theorem~\ref{thm:tailwise-derived} as a removed threshold rather than as an
interior point of a tail.

\begin{proposition}[One-chart equivalence]
\label{prop:app-one-chart}
Let \(a\neq0\) and let \(F\) be a connected component of \(W^{\Gm}\).  Write
\(\mathfrak S^{\,0}_a\) for the stack of rotation-fixed sections over the chart
\(P^0\) with exponent \(a\), imposing membership in \(W^{\mathrm{gen}}\) and
limit in \(F\).  Then evaluation at \(z=1\) is an equivalence of derived stacks
\[
 \operatorname{ev}_1\colon\mathfrak S^{\,0}_a
 \xrightarrow{\ \sim\ }[W_F^a/\Gm],
\]
natural in the simplicial test algebra, compatible with arbitrary derived base
change, and independent of the chosen invariant affine chart.  In particular
\(\mathfrak S^{\,0}_a\) has no derived structure, and its classical truncation
is the Bia{\l}ynicki--Birula stratum \(W_F^a\) modulo the gauge group,
including stabilizers.
\end{proposition}

\begin{proof}
Work first on \(\operatorname{Spec}A\) and before taking the quotient by the
gauge group.  By Lemma~\ref{lem:app-graded-extension} the sections over
\(P^0\) with exponent \(a\) and values in \(\operatorname{Spec}A\) are, via
evaluation at \(1\), the \(R\)-points of \(\operatorname{Spec}A^a\); this is an
isomorphism of functors on simplicial commutative \(\C\)-algebras, hence
commutes with base change, since \eqref{eq:app-graded-map} is levelwise and
does not involve \(R\).  For \(a>0\) the ideal \(I_a\) is generated by the
negative weight spaces, so \(\operatorname{Spec}A^a\) is the locus where
\(\lim_{t\to0}t\cdot x\) exists; for \(a<0\) it is the locus where the opposite
limit exists.  The subalgebra of positive-degree elements of \(A^a\) is an
ideal, so the projection to the weight-zero part is a ring map and defines the
limit morphism \(\operatorname{Spec}A^a\to\operatorname{Spec}A_0\) onto the
fixed locus of the chart.  The preimage of a connected component \(F\) is open
and closed in \(\operatorname{Spec}A^a\); intersecting with the open
\(W^{\mathrm{gen}}\) imposes generic semistability.  This is the chartwise
description of \(W_F^a\).

Chart independence.  The chartwise condition of
Lemma~\ref{lem:app-graded-extension}\textup{(b)} is that the limit exists
\emph{inside the given chart}, and that condition is genuinely chart
dependent.  Take \(W=\mathbf P^1\) with \(t\cdot y=ty\) in the target
coordinate \(y\), and \(a=+1\).  On \(V_0=\operatorname{Spec}\C[y]\) the
coordinate has weight \(+1\), no weight space is killed, and the chart imposes
no condition at all; on \(V_\infty=\operatorname{Spec}\C[y^{-1}]\) the
coordinate \(y^{-1}\) has weight \(-1\), so \(I_a=(y^{-1})\) and the chart
imposes \(y^{-1}=0\).  On the overlap the two conditions disagree at every
point.  Separatedness makes an
extension unique once it exists in a given chart; it does not make it exist.
The gluing instead follows from invariance of the charts.  Let
\(x\in W_F^a\) and let \(y=\lim_{t\to0}t^ax\in F\).  By Sumihiro's theorem
choose a \(\Gm\)-invariant affine open \(U\ni y\).  Since \(U\) is open and
the completed orbit map \(\mathbf A^1\to W\), \(t\mapsto t^ax\), is
continuous, \(t^ax\in U\) for all small \(t\neq0\); since \(U\) is invariant,
\(x\in U\), and then the whole completed orbit lies in \(U\).  Hence
\(W_F^a\) is covered by the invariant affine charts \emph{meeting \(F\)}, and
on each of those the previous paragraph is the correct chartwise description;
charts disjoint from \(F\) compute a different stratum and are excluded.
Uniqueness of the extension then makes the identification
\(\operatorname{ev}_1\) compatible on overlaps, so it globalizes.  The family
form needs one more word, since a section over \(P^0_R\) must factor through a
single chart.  It does: the invariant affine charts meeting \(F\) are opens of
\(W\) covering \(W_F^a\) by the previous paragraph, so for an \(R\)-point
\(x\) of \(W_F^a\) their preimages \(x^{-1}(U)\) are an open cover of
\(\operatorname{Spec}R\); Zariski-locally on \(\operatorname{Spec}R\) the point
\(x\) factors through one such \(U\), and then so does the whole section
\(z\mapsto z^ax\), by invariance of \(U\) again.

Descent to the quotient stack.  A rotation-fixed section over \(P^0_R\) with
values in \([W/\Gm]\) is, locally on \(\operatorname{Spec}R\), a trivialization
of the \(\Gm\)-bundle together with a section as above.  Here we use that a
rotation-equivariant \(\Gm\)-bundle on \(P^0_R\) is locally trivial on
\(\operatorname{Spec}R\): equivariantly, such a bundle is a graded invertible
\(R[z]\)-module, and the graded Picard group of \(R[z]\) is
\(\operatorname{Pic}(R)\times\mathbf Z\), the second factor recording the
clutching exponent.  The nonequivariant analogue is false for nonreduced
\(R\), which is why the equivariant statement is the one used.  Two
trivializations
differ by a rotation-invariant gauge transformation, and the rotation-invariant
gauge transformations over \(P^0\) are the constants \(\Gm(R)\).  Quotienting
by this action turns the previous paragraph into the asserted equivalence with
\([W_F^a/\Gm]\), stabilizers included.

Absence of derived structure.  The chart with its rotation action is the
quotient stack \([\A^1/\C^\times_{\mathrm{rot}}]\).  Its quasicoherent
cohomology vanishes in positive degrees, since cohomology on \(\A^1\) vanishes
there and taking rotation invariants is exact in characteristic zero.  Hence
the deformation theory of \(\mathfrak S^{\,0}_a\) relative to a point is
concentrated in degree zero, so the stack is smooth and its structure sheaf is
discrete.  This is consistent with Bia{\l}ynicki--Birula: for smooth \(W\) the
fixed components and their attracting loci are smooth.
\end{proof}

\begin{proposition}[Two-chart descent]
\label{prop:app-descent}
With \(\mathfrak S_{a_-,a_+}\) as in Definition~\ref{def:app-fixed-section},
evaluation at \(1\) is an equivalence of derived stacks
\[
 \operatorname{ev}_1\colon\mathfrak S_{a_-,a_+}
 \xrightarrow{\ \sim\ }
 \mathfrak Z=
 [W_{F_-}^{a_-}/\Gm]\mathop{\times}^{\mathbf R}_{[W/\Gm]}
 [W_{F_+}^{a_+}/\Gm],
\]
and the inverse equivalence equips a family \(x\) with the two prescribed
clutching bundles and the sections \(z\mapsto z^{a_\pm}x\).  By
Lemma~\ref{lem:app-sign}, the right-hand side depends on the pair
\((a_-,a_+)\) only through the two signs and the two components
\(F_-,F_+\).
\end{proposition}

\begin{proof}
The two charts cover \(P\) and meet in the open orbit.  A section over \(P_R\)
is a pair of sections over the two charts together with an identification of
their restrictions to the open orbit, and mapping into a stack makes this a
homotopy limit rather than a strict equality of restrictions; that homotopy
limit is the homotopy fibre product of the two chart-section stacks over the
section stack of the open orbit.  The rotation-fixed sections over the open
orbit are, by evaluation at \(1\) and the argument of
Proposition~\ref{prop:app-one-chart} with \(I_a=0\), the stack \([W/\Gm]\); the
restriction maps from the two charts are the two stratum inclusions.  Applying
Proposition~\ref{prop:app-one-chart} to each chart gives the displayed
equivalence.  Since evaluation at \(1\) is the common value of the two
extensions, the inverse construction reads off \(x\) and reconstructs the two
sections by \eqref{eq:app-graded-map}, which is where the exponents reappear.
\end{proof}

\subsection{Tangent complexes and the
\texorpdfstring{\(2\)-commutative}{2-commutative} square}
\label{app:square}

All complexes in this subsection are relative to the fixed moduli \(\mathfrak
M\) of marked domains.  Let \(\pi\colon\mathcal P\to\mathfrak Z\) be the
universal domain, \(u_a\colon\mathcal P\to[W/\Gm]\) the universal gauged map
with exponents \(a=(a_-,a_+)\), and
\[
 V:=\operatorname{ev}_1^*T([W/\Gm])
\]
the pullback of the tangent complex of the quotient stack along evaluation at
\(1\).  Define \(F^0_{\leq0}V\) and \(F^\infty_{\leq0}V\) to be the subcomplexes
of \(V\) consisting of the rotation-invariant sections of
\(u_a^*T([W/\Gm])\) over the chart at \(0\), respectively at \(\infty\),
viewed inside \(V\) by restriction to the open orbit.  These are the extension
subcomplexes of the in-text display; the subscript records that they select the
tangent weights \(w\) with \(aw\leq0\) at the corresponding limit point, as
computed in Lemma~\ref{lem:app-cech}, and only the intrinsic description just
given is used below.  No equivariant splitting of
\(V\) is chosen.

\begin{lemma}[Invariant \v{C}ech presentation]
\label{lem:app-cech}
There are canonical equivalences, natural in \(\mathfrak Z\),
\begin{equation}\label{eq:app-cech}
 \bigl(R\pi_*u_a^*T([W/\Gm])\bigr)^{\mathrm{rot}}
 \simeq
 \bigl[F^0_{\leq0}V\oplus F^\infty_{\leq0}V\longrightarrow V\bigr]
 \simeq
 T_{\mathfrak Z/\mathfrak M},
\end{equation}
where the middle complex sits in degrees \(0\) and \(1\) and its differential
is the difference of the two restrictions to the open orbit.  Moreover
\(F^0_{\leq0}V\) and \(F^\infty_{\leq0}V\) are the pullbacks of
\(T([W_{F_-}^{a_-}/\Gm])\) and \(T([W_{F_+}^{a_+}/\Gm])\), and all three terms
depend on \(a\) only through the two signs.
\end{lemma}

\begin{proof}
The pullback \(u_a^*T([W/\Gm])\) is a two-term complex of vector bundles on
\(\mathcal P\), and the two charts are affine over \(\mathfrak Z\), so the
\v{C}ech complex of the two-chart cover computes \(R\pi_*\).  It is
\([\,\Gamma(P^0)\oplus\Gamma(P^\infty)\to\Gamma(\Gm)\,]\) with the difference
differential.  Taking \(\C^\times_{\mathrm{rot}}\)-invariants is exact in
characteristic zero, because the rotation group is linearly reductive and
invariants are the weight-zero summand of a weight decomposition; so invariants
may be taken termwise.  By definition the invariants of the three terms are
\(F^0_{\leq0}V\), \(F^\infty_{\leq0}V\) and \(V\), the last because a
rotation-invariant section over the open orbit is determined by its value at
\(z=1\).  This gives the first equivalence.

For the identification of the two extension subcomplexes, apply
Lemma~\ref{lem:app-graded-extension} to the module of sections rather than to
the ring, with one change of sign.  A vector field \(D\) has weight \(w\) when
\(D(A_v)\subset A_{v+w}\), so tangent weights are opposite to the weights of
the coordinate functions, and along the universal section a tangent vector of
weight \(w\) acquires \(z^{-aw}\) rather than \(z^{aw}\).  Hence on an
invariant affine chart a graded module \(M=\bigoplus_wM_w\) of sections gives
\((M\otimes_AR[z])^{\mathrm{rot}}=\bigoplus_{aw\leq0}(M_w\otimes_AR)\) and
\((M\otimes_AR[z^{\pm1}])^{\mathrm{rot}}=\bigoplus_w(M_w\otimes_AR)\).  This
nonpositive half is what the subscript in \(F^0_{\leq0}V\) and
\(F^\infty_{\leq0}V\) records.  For \(W=\A^1\) with \(t\cdot y=t^ky\) and
\(a>0\), the attracting locus is all of \(\A^1\) and its tangent direction has
weight \(-k\), which the displayed condition retains.  Taking
\(M\) to be the module of sections of \(T([W/\Gm])\) identifies the invariants
over the chart with the tangent complex of the attractor subscheme, which is
the content of Bia{\l}ynicki--Birula's description of the attracting locus of a
smooth \(\Gm\)-variety.  The same statement in family form over a simplicial
base is the module version of Proposition~\ref{prop:app-one-chart}.

The second equivalence in \eqref{eq:app-cech} is the tangent complex of a
homotopy fibre product: for \(\mathfrak Z\) as in
Proposition~\ref{prop:app-descent} there is a fibre sequence
\(T_{\mathfrak Z/\mathfrak M}\to T_{-}\oplus T_{+}\to V\) with
\(T_\pm=T([W_{F_\pm}^{a_\pm}/\Gm])\) pulled back, which is the displayed
two-term complex.  Finally, sign dependence follows from
Lemma~\ref{lem:app-sign} applied to the module computation above.
\end{proof}

\begin{remark}[Where the degree-one term comes from]
\label{rem:app-degree-one}
The complex \eqref{eq:app-cech} has \(h^1\) equal to the cokernel of
\(F^0_{\leq0}V\oplus F^\infty_{\leq0}V\to V\), which is nonzero exactly when
the two attracting loci fail to meet transversely at the point in question.
The gauge direction contributes to all three terms and cancels in the
difference; on the generic semistable locus the stabilizers are finite, so
\(h^{-1}\) vanishes and the complex is supported in degrees \(0\) and \(1\).
Along one tail the two signs are constant and the zero-weight changes were
removed as thresholds, so the two subcomplexes and the restriction map are
constant, while the universal section \(z\mapsto z^{a_\pm}x\) is not.
\end{remark}

\begin{lemma}[Woodward's obstruction theory as a truncation]
\label{lem:app-truncation}
Put
\[
 E_a=\bigl((R\pi_*u_a^*T([W/\Gm]))^{\mathrm{rot}}\bigr)^\vee.
\]
The cotangent complex of the derived mapping stack gives
\(E_a\simeq\tau_{[-1,0]}L_{\mathfrak S_{a_-,a_+}/\mathfrak M}\), and under this
identification the morphism
\(E_a\to L_{t_0\mathfrak S_{a_-,a_+}/\mathfrak M}\) of
\cite[Definition~7.13]{WoodwardQKIII}, in its relative form over the twisted
domain stack \cite[Section~8.3]{WoodwardQKIII},
\cite[Example~6.6(c)]{WoodwardQKII}, is the composite
\[
 \tau_{[-1,0]}L_{\mathfrak S_{a_-,a_+}/\mathfrak M}
 \longrightarrow L_{\mathfrak S_{a_-,a_+}/\mathfrak M}
 \longrightarrow L_{t_0\mathfrak S_{a_-,a_+}/\mathfrak M}.
\]
\end{lemma}

\begin{proof}
The first assertion is the cotangent complex formula for a mapping stack
\cite{ToenVezzosiHAGII}: the
relative cotangent complex of \(\mathfrak S_{a_-,a_+}\) over \(\mathfrak M\) is
the pushforward along \(\pi\) of the pullback of the cotangent complex of the
target, dualized as above, and the rotation-invariant part appears because the
sections are rotation-fixed.  Both sides are perfect of amplitude
\([-1,0]\) by Lemma~\ref{lem:app-cech}.

For the second assertion, both morphisms are the map classifying the
deformation theory of the same universal object: Woodward's is obtained from
the Kodaira--Spencer map of the universal gauged map relative to the domain
stack, which is exactly the map that the derived structure of the mapping stack
induces on the classical truncation.  In the presentation of
Lemma~\ref{lem:app-cech} both are represented by the same map, namely the dual
of the map that sends a first-order deformation of a section to the induced
pair of deformations over the two charts.  We use here the standard comparison
between a classical obstruction theory built from the deformation theory of a
universal object and the truncation of the derived cotangent complex; that
comparison is imported from \cite{SchurgToenVezzosi}, not reproved.
\end{proof}

\begin{proposition}[The \(2\)-commutative square]
\label{prop:app-square}
Under the equivalence
\(\operatorname{ev}_1\colon\mathfrak S_{a_-,a_+}\simeq\mathfrak Z\) of
Proposition~\ref{prop:app-descent}, the square
\[
\begin{array}{ccc}
 E_a & \longrightarrow & L_{t_0\mathfrak S_{a_-,a_+}/\mathfrak M}\\
 \big\downarrow\scriptstyle\sim && \big\downarrow\scriptstyle\sim\\
 \tau_{[-1,0]}L_{\mathfrak Z/\mathfrak M}
   & \longrightarrow & L_{t_0\mathfrak Z/\mathfrak M}
\end{array}
\]
is \(2\)-commutative, its left vertical arrow is a quasi-isomorphism, and the two
perfect obstruction theories are isomorphic as morphisms to the cotangent
complex.  Passing to classical truncations, the derived data recover Woodward's
relative perfect obstruction theory, hence his relative virtual class.
\end{proposition}

\begin{proof}
By Lemma~\ref{lem:app-truncation} the upper horizontal map is the truncation
map of \(L_{\mathfrak S_{a_-,a_+}/\mathfrak M}\); the lower horizontal map is
by definition the truncation map of \(L_{\mathfrak Z/\mathfrak M}\).  The
assignment \(\mathfrak X\mapsto\bigl(\tau_{[-1,0]}L_{\mathfrak X/\mathfrak
M}\to L_{t_0\mathfrak X/\mathfrak M}\bigr)\) is functorial for equivalences of
derived stacks over \(\mathfrak M\), and the two vertical arrows are the values
of that functor on \(\operatorname{ev}_1\).  The square is therefore the
naturality square of a natural transformation evaluated at an equivalence, and
the required homotopy is the one supplied by \(\operatorname{ev}_1\).

It remains to check that this homotopy is the one asserted in the proof of
Theorem~\ref{thm:tailwise-derived}, namely the differential of evaluation at
\(1\) on the two-chart universal section.  In the presentation
\eqref{eq:app-cech} this is an equality of representatives rather than a
homotopy.  Indeed, both composites in the square are, after dualizing,
computed on a first-order deformation as follows.  Restrict the deformation of
the section to the two charts; each restriction is a section of the extension
subcomplex at that chart by Lemma~\ref{lem:app-cech}; then restrict both to the
open orbit and take the difference in \(V\).  For the upper route this is the
\v{C}ech differential of \(R\pi_*u_a^*T([W/\Gm])\); for the lower route it is
the structure map of the homotopy fibre product \(\mathfrak Z\).  Both are the
map \(F^0_{\leq0}V\oplus F^\infty_{\leq0}V\to V\) given by the difference of
the two restrictions to the open orbit, and these restrictions are the same
morphisms of sheaves in the two cases, because the universal section restricted
to a chart is the composite of the stratum inclusion with the chart section
\eqref{eq:app-graded-map}, and the latter is an isomorphism onto its image
after evaluation at \(1\).  The two representatives are equal, so the
comparison homotopy may be taken to be the identity.

Since \(\operatorname{ev}_1\) is an equivalence, its induced map on
\(\tau_{[-1,0]}L_{-/\mathfrak M}\) is a quasi-isomorphism; this is the left
vertical arrow.  A commuting square whose verticals are quasi-isomorphisms is
an isomorphism of the two morphisms to the cotangent complex, which is the
assertion for obstruction theories.  Taking \(t_0\) and using
Lemma~\ref{lem:app-truncation} identifies the upper row with Woodward's
relative perfect obstruction theory, and an isomorphism of relative perfect
obstruction theories over a common base induces equality of the associated
relative virtual classes in Vistoli's bivariant Chow groups
\cite[Example~6.6(c)]{WoodwardQKII}.
\end{proof}

\subsection{Orbifold charts, \texorpdfstring{\(\mu_k\)}{mu-k}-equivariance, and
rigidification}
\label{app:rigidification}

Woodward's clutching construction has an orbifold form
\cite[Definition~9.7]{WoodwardQKIII}: one passes to a \(k\)-fold cover of the
open orbit, the cocharacters \(\widetilde\varphi_\pm\) are taken with values in
\(\tfrac1k\mathbf Z\), and the section on a chart is
\(z\mapsto\widetilde\varphi_\pm(z^{1/k})x\).  The chart is then
\([\operatorname{Spec}R[z^{1/k}]/\mu_k]\) with \(\mu_k\) acting by
\(z^{1/k}\mapsto\theta z^{1/k}\) for a primitive \(k\)-th root of unity
\(\theta\).

\begin{proposition}[Equivariance and descent]
\label{prop:app-mu-k}
Write \(a=b/k\) with \(b\in\mathbf Z\).  Then:
\begin{enumerate}[label=\textup{(\alph*)}]
\item Lemma~\ref{lem:app-graded-extension} holds verbatim with \(R[z]\)
replaced by \(R[z^{1/k}]\) and \(aw\in\tfrac1k\mathbf Z\); the attractor ideal
is unchanged, being \((A_w:bw<0)\).
\item The stabilizer element of the clutching datum is
\(\widetilde\varphi(\theta)=\theta^{\,b}\), so the inertia label, the
\(\mu_k\)-weights of the universal bundle at the orbifold point, and the
twisted-sector evaluation depend on \(b\) only through its class modulo \(k\).
Within one tail, consecutive affine degrees differ by \(k\) and the sign of
\(b\) is constant.
\item Taking \(\mu_k\)-invariants is exact in characteristic zero, and the
equivalence of Proposition~\ref{prop:app-descent} and the square of
Proposition~\ref{prop:app-square} are \(\mu_k\)-equivariant.
\item Rigidification replaces the inertia stack by its quotient by the finite
central \(B\mu_k\), \(\bar I_{X,r}=I_{X,r}/B\mu_r\)
\cite[Section~4.3, Proposition~4.3(f),(g)]{WoodwardQKII}.  Pullback along this
gerbe is fully faithful with \(\pi_*\pi^*\simeq\id\), hence conservative, so
the equivalence and the obstruction theory descend.  Stabilizer and
graph-automorphism groups, and their localization multiplicities, are
unchanged.
\end{enumerate}
\end{proposition}

\begin{proof}
(a)  The proof of Lemma~\ref{lem:app-graded-extension} used only that the
target is a graded polynomial extension of \(R\) whose degrees are bounded
below by zero.  With \(z^{1/k}\) of degree \(1/k\) the same splitting applies,
and \(aw<0\) is the condition \(bw<0\).

(b)  The stabilizer of the orbifold point acts through
\(\widetilde\varphi(\theta)\), which for \(G=\Gm\) and
\(\widetilde\varphi(t)=t^{b/k}\) is \(\theta^{\,b}\); this depends only on
\(b\bmod k\).  The isotropy representation of the universal bundle at the
orbifold point is determined by the same element, and the twisted-sector
evaluation of \cite[Corollary~9.10]{WoodwardQKIII} is the corresponding
component of the inertia stack.  Fixing the stabilizer congruence class of
Theorem~\ref{thm:tailwise-derived} is exactly fixing \(b\bmod k\), so a tail
consists of the \(b\) in one class with one sign, and its consecutive elements
differ by \(k\).

(c)  Exactness of invariants under a finite group in characteristic zero is
standard, and it is the only property of \(\mu_k\) used in
Lemma~\ref{lem:app-cech}.  All the constructions of
Sections~\ref{app:attractor} and~\ref{app:square} are defined by the graded
formula \eqref{eq:app-graded-map} and by restriction morphisms, both of which
commute with the \(\mu_k\)-action on \(z^{1/k}\); the equivalence and the
square are therefore diagrams of \(\mu_k\)-equivariant objects and morphisms.

(d)  A \(B\mu_k\)-gerbe in characteristic zero has \(\pi_*\pi^*\simeq\id\) on
quasicoherent complexes, because \(\pi_*\) takes the weight-zero part of the
\(\mu_k\)-decomposition and a pulled-back complex is concentrated in weight
zero.  Hence \(\pi^*\) is fully faithful and conservative, and a morphism of
pulled-back objects which is \(\mu_k\)-equivariant descends and is an
equivalence downstairs precisely when it is one upstairs.  The equivalence of
Proposition~\ref{prop:app-descent} and the square of
Proposition~\ref{prop:app-square} are of that form by (c).  The fibre products
of \cite[equations~(54) and~(57)]{WoodwardQKIII} are taken over the
unrigidified inertia stack, so the constructions above take place there;
passage to the rigidified inertia stack used for the evaluation maps changes
rational cohomology only by the factors of \(r\) on \(r\)-twisted sectors
recorded in \cite[Proposition~4.3(f),(g)]{WoodwardQKII}, and those factors
depend only on the order of the stabilizer, which is constant along a tail by
(b).  The parametrized principal component has no domain automorphisms; the
fixed bubble factors retain theirs, and their automorphism groups are the ones
appearing in the localization multiplicities.
\end{proof}

\subsection{Attached bubbles, cutting, and reduction between Artin levels}
\label{app:cutting}

\begin{proposition}[Attached factors and cutting]
\label{prop:app-cutting}
Let \(\mathfrak F\) be the fixed graph stack obtained from
\(\mathfrak S_{a_-,a_+}\) by attaching the fixed stable-map factors of
\cite[equations~(54)--(57)]{WoodwardQKIII} along the two inertia evaluations,
as a derived fibre product.  Then:
\begin{enumerate}[label=\textup{(\alph*)}]
\item the relative cotangent complex of \(\mathfrak F\) fits in the exact
triangle whose outer terms are the direct sum of the factor obstruction
theories and the pullback of the cotangent complex of the inertia stack;
\item the identification of Proposition~\ref{prop:app-square} is the identity
on the stable-map and inertia terms, hence induces an isomorphism of these
triangles;
\item consequently the identification is compatible with the cutting axiom
\cite[Proposition~7.14(b)]{WoodwardQKIII}, whose proof identifies the gauged
moduli stack with a fibre product over the diagonal
\cite[Proposition~5.21]{WoodwardQKII}, and it identifies the resulting virtual
classes;
\item every construction above is a derived fibre product along finitely many
morphisms of stacks over \(\C\); the coefficient ring enters only through the
localization coefficients, so the identifications commute with the reduction
maps between finite Artin levels.
\end{enumerate}
For the irreducible unstable convention there is no stable-map factor: by
\cite[Example~9.15]{WoodwardQKIII} the domain is irreducible and the normal
complex is the moving index alone, so the node and attaching factors of
\eqref{eq:virtual-normal-euler} are absent and the square of
Proposition~\ref{prop:app-square} applies with no gluing term.
\end{proposition}

\begin{proof}
(a) is the cotangent triangle of a derived fibre product, applied to
\(\mathfrak F=\mathfrak F_-\times^{\mathbf R}_{I}\mathfrak F_+\) with \(I\) the
inertia stack of the quotient: the triangle reads
\(\operatorname{ev}^*L_I\to L_{\mathfrak F_-}\oplus L_{\mathfrak F_+}\to
L_{\mathfrak F}\) after pullback to \(\mathfrak F\).  This is the derived form
of the exact triangle by which Woodward derives cutting-edge compatibility for
the index complexes \cite[Equation~(35) and Section~7.1]{WoodwardQKIII}.

(b)  The degree shift changes only the principal factor.  The attached
stable-map factors and their obstruction theories are the fixed proper
stable-map data at a fixed ordinary degree and marking set, and the inertia
term is determined by the stabilizer congruence class, which is fixed along a
tail by Proposition~\ref{prop:app-mu-k}(b).  Hence the comparison map is the
identity on those two terms, and Proposition~\ref{prop:app-square} supplies the
map on the third; a map of triangles with two identity legs and one
quasi-isomorphism is an isomorphism of triangles.

(c)  Woodward's cutting axiom expresses the virtual class of the glued type as
the Gysin pullback \(\Delta^!\) of the virtual class of the cut type, using the
identification of the moduli stack with a fibre product over the diagonal.  The
Gysin map is determined by the obstruction theory of the factors and the
regular embedding \(\Delta\); by (a) and (b) both are identified, so the two
cut classes agree.

(d)  Each construction is one of: the equivalence of
Proposition~\ref{prop:app-descent}, a homotopy fibre product along evaluation
maps, or the truncation map of a cotangent complex.  All are morphisms of
stacks and of complexes over \(\C\), independent of the Novikov and bulk
coefficients; a finite Artin quotient enters only when the resulting virtual
classes are paired with localization coefficients.  The identifications are
therefore compatible with every reduction \(R_{N+1}\to R_N\).  Compatibility
here is a statement about the identifications, not about the numbers they
produce: the identified stacks, obstruction theories, and virtual classes carry
no \(N\), so there is nothing for a reduction to act on, and the assertion is
automatic rather than a level-by-level verification.  The paired
localization coefficients do depend on \(N\), and for them compatibility means
the usual thing, that reduction sends the level-\(N+1\) pairing to the
level-\(N\) one.
\end{proof}

\begin{remark}[Imported statements]
\label{rem:app-imports}
Two comparisons are used above without proof.  The first, from
\cite{ToenVezzosiHAGII}, is the cotangent complex formula for a derived
mapping stack, used in Lemma~\ref{lem:app-truncation} to identify \(E_a\)
with \(\tau_{[-1,0]}L_{\mathfrak S_{a_-,a_+}/\mathfrak M}\).  The second,
from \cite{SchurgToenVezzosi}, is the identification of a classical perfect
obstruction theory built from the deformation theory of a universal object
with the classical truncation of the derived one; it is used in the same
lemma to recognize \cite[Definition~7.13]{WoodwardQKIII} in that
truncation.  Everything else is
either the graded computation of Lemma~\ref{lem:app-graded-extension} and its
module form, or a statement of Woodward's cited at the locators given at the
start of this appendix.  In particular, no derived statement is used at any
point where a virtual class or a localization contribution is finally
evaluated: those are read off from the classical truncations, in the form in
which \cite{WoodwardQKIII} states them.
\end{remark}
